% !PRINTED 566
\weberpage{566}
%<<CONT
\[ \sum^{s} \varepsilon^{(\lambda+\mu)\,\mathrm{ind}\,s(t+1)} r^{s(t+1)} = \sum^{s} \varepsilon^{(\lambda+\mu)\,\mathrm{ind}\,s} r^s = (\varepsilon^{\lambda+\mu}, r), \]
and formula (4) yields
\weq{5}{\frac{(\varepsilon^\lambda, r)(\varepsilon^\mu, r)}{(\varepsilon^{\lambda+\mu}, r)} = \psi_{\lambda,\mu}(\varepsilon),}
if for abbreviation
\weq{6}{\sum_{1,\,n-2}^{t} \varepsilon^{\mu\,\mathrm{ind}\,t - (\lambda+\mu)\,\mathrm{ind}(t+1)} = \psi_{\lambda,\mu}(\varepsilon)}
is set, and $\lambda + \mu$ is assumed not divisible by $n - 1 = m$.
These functions $\psi$ are then composed of the $(n-1)$th roots of unity
$\varepsilon$ with integral coefficients, and are therefore
known if the $(n-1)$th roots of unity are assumed to be known.
They are numbers of the field $R(\varepsilon)$.

We now assume $\lambda$ divisible by $\mu$, and accordingly replace $\lambda$
by $\mu\lambda$ and set
\weq{7}{\alpha = \varepsilon^\mu,}
so that $\alpha$ may also be a non-primitive $(n-1)$th root of unity,
e.g.\ an $e$th root of unity (if, as above, $m = ef$
and $\mu$ is divisible by $f$, say taken $= f$), where
however the case $\alpha = 1$, that is $e = 1$, is to be excluded.
Then we set
\weq{8}{\sum_{1,\,n-2}^{t} \alpha^{\mathrm{ind}\,t - (\lambda+1)\,\mathrm{ind}(t+1)} = \psi_\lambda(\alpha);}
here $\psi_\lambda(\alpha)$ is a number of the field $R(\alpha)$, which can easily
be computed with the aid of an index table; $\lambda$ can here be taken
with respect to the modulus $e$. Formula (5) gives for this
assumption
\weq{9}{(\alpha^\lambda, r)(\alpha, r) = (\alpha^{\lambda+1}, r)\psi_\lambda(\alpha),}
and this formula holds so long as $\mu(\lambda+1)$ is not divisible by $n - 1$;
thus, if $\mu$ is coprime to $m$, and hence $\alpha$ is a
primitive $m$th root of unity, for $\lambda = 1, 2 \ldots n-3$; and
if $f$ is the greatest common divisor of $\mu$ and $m$ and
$m = ef$, for $\lambda = 1, 2 \ldots e-2$.
Formula (3) now gives, if $\alpha^\lambda$ is not $= 1$,
\weq{10}{(\alpha^\lambda, r)(\alpha^{-\lambda}, r) = (-1)^{\mu\lambda} n.}
Formulae (9), (10) suffice to determine not only the $r$ themselves,
but also all the periods $\eta$, by radicals
whose degrees are the primes dividing $m$,
and among which the functions $\psi_\lambda$ occur as the known
quantities.
